\documentclass[12pt]{article}
\usepackage[utf8]{inputenc}
\usepackage[T1]{fontenc}
\usepackage[margin=1in]{geometry}
\usepackage{setspace}
\usepackage{xcolor}
\usepackage[colorlinks=true, urlcolor=blue!60!black, linkcolor=blue!60!black]{hyperref}
\setlength{\parindent}{2em}
\setlength{\parskip}{1.25em}
\renewcommand{\baselinestretch}{1.0}
\begin{document}
\begin{center}
{\LARGE \textit{Arizona v. United States}, 567 U.S. 387 (2012)\par}
\vspace{0.8em}
{\large Case Brief --- Four-Page Edition\par}
\end{center}
\vspace{1em}

\section{Memory Jogger}

The Court struck down three of four challenged provisions of Arizona's S.B.~1070 as preempted by federal immigration law---state crimes for alien-registration failures and unauthorized work, and state warrantless arrests for removability---while declining to facially enjoin the ``show me your papers'' status-verification mandate before state courts could construe it (\href{https://supreme.justia.com/cases/federal/us/567/387/}{Justia}).

\section{Detailed Case Facts}

In April 2010 Arizona enacted S.B.~1070 to ``discourage and deter the unlawful entry and presence of aliens and economic activity by persons unlawfully present in the United States,'' adopting an official state policy of ``attrition through enforcement.'' The United States brought a pre-enforcement facial challenge under the Supremacy Clause against four provisions: \textbf{Section 3} made the ``willful failure to complete or carry an alien registration document'' a state misdemeanor; \textbf{Section 5(C)} made it a state crime for an unauthorized alien to seek or perform work, although IRCA penalizes employers and imposes only civil consequences on unauthorized employees; \textbf{Section 6} authorized warrantless state arrests on probable cause that a person had committed a removable offense; and \textbf{Section 2(B)} required officers making a lawful stop, detention, or arrest to make a ``reasonable attempt'' to determine immigration status with ICE when reasonable suspicion of unlawful presence exists. The United States challenged the law solely on preemption grounds, not on Fourth Amendment or racial-profiling claims.

\section{Procedural History}

The district court preliminarily enjoined all four provisions. A divided Ninth Circuit panel affirmed, with Judge Bea dissenting as to Sections 2(B) and 6. Arizona petitioned directly for certiorari, which was granted in December 2011; the Court decided the case on June 25, 2012.

\section{Judicial Votes}

Affirmed in part, reversed in part, and remanded: \textbf{5--3}, with Justice Kagan recused.
\begin{itemize}
\item \textbf{Majority:} Justice Kennedy, joined by Chief Justice Roberts and Justices Ginsburg, Breyer, and Sotomayor.
\item \textbf{Concurring in part and dissenting in part:} Justice Scalia (would uphold all four provisions).
\item \textbf{Concurring in part and dissenting in part:} Justice Thomas (would uphold all four on textualist grounds).
\item \textbf{Concurring in part and dissenting in part:} Justice Alito (agreed on Sections 2(B) and 3; would uphold Sections 5(C) and 6).
\end{itemize}

\section{Holding}

Federal immigration law displaces state statutes that intrude into a field Congress occupies exclusively, such as alien registration, even when the state law mirrors federal standards, or that stand as an obstacle to Congress's deliberate balance and federal enforcement discretion, such as criminalizing unauthorized work or granting state officers unilateral warrantless arrest power for civil removability. A mandate to verify status with ICE during lawful stops is not facially preempted before state-court construction, provided it does not prolong detention solely to check status.

\section{Analysis of Opinions}

\textbf{Majority (Kennedy).} Federal power rests on the Naturalization Clause and inherent sovereignty over foreign relations: foreign nations must deal with one national sovereign, not fifty States. \emph{Section 3}: under \textit{Hines v. Davidowitz}, Congress has occupied the field of alien registration through a ``single integrated and all-embracing system,'' so even parallel state penalties are barred. \emph{Section 5(C)}: IRCA made a deliberate choice to sanction employers, not employees, criminally; a ``conflict in technique'' can disrupt that calibrated scheme. \emph{Section 6}: unlawful presence is generally civil, and the provision gave state police broader arrest authority than Congress gave trained federal officers. \emph{Section 2(B)}: Congress encouraged federal-state information sharing, and state courts could construe the provision narrowly, so a facial injunction was premature; as-applied challenges remain open.

\textbf{Scalia.} States possess an inherent sovereign power to exclude unwanted noncitizens unless Congress unequivocally abrogates it. The challenged provisions merely enforce federal restrictions, and the executive's nonenforcement policies (he cited the 2012 DACA memorandum) cannot preempt state laws that faithfully enforce what Congress enacted.

\textbf{Thomas.} The Supremacy Clause gives preemptive force only to the ordinary meaning of enacted federal text, not to implied purposes; finding no textual conflict, he would uphold every provision.

\textbf{Alito.} He agreed on Sections 2(B) and 3, calling the claim that shifting enforcement priorities can preempt state law an ``astounding assertion of federal executive power,'' but faulted the majority for giving ``short shrift'' to the presumption against preemption on Section 5(C) and would have sustained Section 6 against a facial challenge.

\section{Examples: Future Applications}

\begin{enumerate}
\item \textbf{Same-Side (Preempted): State Crime of Unlawful Entry.} A border State criminalizes unauthorized entry and lets state judges order noncitizens to return abroad. Field and obstacle preemption both apply: unilateral state removal orders usurp federal foreign-relations authority and prosecutorial discretion.
\item \textbf{Same-Side (Preempted): Rental Ban.} A city fines undocumented individuals who occupy rental housing. Like Sections 3 and 5(C), it imposes local sanctions to drive out noncitizens whom federal authorities may choose not to remove.
\item \textbf{Opposite-Side (Not Preempted): Tax-Form Identity Theft.} A neutral identity-theft law punishing use of another's Social Security number on W-4 forms applies to everyone and does not target unauthorized work, as \textit{Kansas v. Garcia} (2020) later held.
\item \textbf{Opposite-Side (Not Preempted): Licensing and E-Verify.} Revoking licenses of employers who knowingly hire unauthorized workers falls within IRCA's licensing savings clause, as \textit{Chamber of Commerce v. Whiting} held.
\item \textbf{Fence-Sitter: 48-Hour ICE Detainers.} A State requires jails to hold inmates up to 48 hours past release on ICE request. Responding to a federal request may be ``cooperation,'' but \textit{Arizona} warned that holding someone past release solely for civil immigration purposes raises serious Fourth Amendment and preemption concerns.
\end{enumerate}

\section{Critique}

\textbf{Civil rights.} Lucas Guttentag argued that by divorcing preemption from antidiscrimination norms the Court obscured the civil-rights harms of ``attrition through enforcement'' laws; litigation over Section 2(B) led Arizona in 2016 to settle, sharply restricting prolonged stops and race-based immigration inquiries.

\textbf{Text and federalism.} Originalist and textualist critics argue that inferring preemption from congressional silence or failed bills undermines the Supremacy Clause, which makes enacted ``Laws'' supreme. Critics also note a tension: the Court rejected executive priorities as a basis for preempting Section 2(B) yet elevated executive discretion into a de facto preemption shield for Sections 3 and 6.

\section{Key Quotations}

\begin{quote}
``The Government of the United States has broad, undoubted power over the subject of immigration and the status of aliens.''
\end{quote}
\noindent --- Justice Kennedy, for the Court, 567 U.S. at 394.

\begin{quote}
``Where Congress occupies an entire field, as it has in the field of alien registration, even complementary state regulation is impermissible. Field preemption reflects a congressional decision to foreclose any state regulation in the area, even if it is parallel to federal standards.''
\end{quote}
\noindent --- Justice Kennedy, for the Court, 567 U.S. at 401.

\begin{quote}
``Under the Supremacy Clause, pre-emptive effect is to be given to congressionally enacted laws, not to judicially divined legislative purposes.''
\end{quote}
\noindent --- Justice Thomas, concurring in part and dissenting in part, 567 U.S. at 440.

\end{document}
